\section{Background and Related Work}

\subsection{Layout targets, representations and measurement}
Panoramic layout estimation represents indoor enclosures through wall turns and their ceiling and floor boundaries \cite{c2ref01,horizon2019}. MatterportLayout and ZInD supply references across diverse indoor scenes \cite{zou2021manhattan,cruz2021zind}, but choosing the enclosed space remains part of the annotation task. Bi-Layout distinguishes layouts that stop at openings from layouts that extend into visible adjoining spaces \cite{c2ref06}. Recent efforts have also scrutinized reference fidelity: ReplicaPano combines and refines annotations across viewpoints of the same room to construct a consistent layout \cite{dong2024panocontext}, while ToF-360 uses building geometry to annotate occluded elements and documents a missing wall plane in a MatterportLayout example \cite{kanayama2025tof360}. These studies show why a reliable reference depends on both the intended enclosure and the evidence used to trace it. A camera beneath a shared doorframe makes this distinction particularly consequential: either adjoining room or their combined space may be selected, with different closure boundaries. The doorway-boundary setting makes these target choices explicit before individual and collective layouts are compared.

Layout evaluation uses overlap, corner, pixel and depth measures under different representations and benchmark protocols \cite{horizon2019,horizoneval,zou2021manhattan}. For freely marked upper endpoints, reducing the layout to a mean-height prism can conceal opposing local height errors. Likewise, omitting a shallow recess may substantially change the wall sequence while barely changing floor overlap. These cases motivate measuring the geometry behind a score. The proposed $Q$ combines extent, boundary placement, orthogonal regularity and global and local height measurements, retaining interpretable components that can be examined alongside structural differences.

\subsection{Annotator differences and spatial aggregation}
Repeated-labeling models account for worker-specific error rates, multidimensional expertise and item difficulty \cite{c2ref13,c2ref08,c2ref09}, and redundant annotations can improve label accuracy \cite{c2ref19,c2ref25}. Other approaches examine the structure of disagreement itself. Studies of divergent interpretations identify annotator groups with recurring labeling patterns \cite{c2ref07}, while CrowdTruth relates disagreement to workers, items and labels \cite{c2ref11}. Research on textual inference and human label variation likewise treats differing judgments as information about the task and its interpretations \cite{c2ref10,plank2022variation}. In panoramic annotation, such differences appear both in which wall turns people retain and in where they place shared corners. Repeated simplification of recesses or retention of short walls therefore provides information about individual annotation patterns as well as geometric accuracy.

Spatial aggregation methods differ in the units they combine. STAPLE estimates segmentation truth and annotator performance from pixel labels \cite{c2ref14}, whereas boundary-based segmentation aggregation forms planar cells from overlapping submitted polygons \cite{lee2018}. A general annotation-aggregation model combines distance-based selection, decomposition, merging and object matching across simple, complex and multi-object tasks \cite{c2ref17}. For panoramic layouts, the aggregation unit must also preserve the cyclic wall sequence and the relationship between upper and lower endpoints, even when annotators mark different numbers of turns. Paired-corner structural consensus combines these observations through cyclic matching, distinct-person support and periodic position aggregation. This representation makes it possible to trace both supported turns and their upper boundaries into the output, complementing the floor coverage captured by BEV region voting.

The choice of aggregation unit also makes member contributions observable. Replacing an annotator can add support for a missed turn, weaken support for a recess or shift a shared boundary, distinctions that average worker or class quality alone cannot describe. Separating class proportions from specific membership, and comparing the same exchanged members against different retained teammates, reveals how individual patterns interact in the collective layout.

\subsection{Task conditions and cross-view histories}
Interactive layout workflows combine proposed geometry with human editing. PanoAnnotator supports initialization, refinement and geometric constraints \cite{pano2018}; ZInD combines automated layout proposals and floor-plan assembly with human correction and verification \cite{cruz2021zind}. In clinical text annotation, fully pre-populated suggestions reduced experts' initiative to create additional annotations \cite{levy2021suggestions}. Annotation guidelines and familiarity with the data also shape task performance \cite{c2ref21,c2ref12}. These findings motivate distinguishing the presence of an edit from the correction of a defect. Manual and Semi submissions provide complementary observations of layouts drawn from scratch and responses to initialized geometry.

Multi-view layout methods use complementary observations of the same room: PSMNet merges layout information across viewpoints, and 360-MLC uses cross-view consistency for self-training and parameter selection \cite{c2ref23,c2ref24}. Shared room geometry offers a corresponding opportunity to study human annotation: changing the viewpoint changes visibility while retaining the room being described. The central questions are whether an annotator repeatedly retains or simplifies structure, which local corner choices change, and whether a specified member replacement has a similar effect across views. Structural histories and quality histories provide distinct evidence about what can transfer between these observations.
